\subsection{Minimal reproducibility objects and admission order}
The necessary deposit has three linked resources and a reconstruction script, not only a paper PDF. The first is a \emph{36-row official-source register}: source ID, layer, country, canonical document title, issuing authority, policy/status, precise issue/effective date and edition, controlling language, direct URL, quoted clause and page, source capture date, source file hash, E/P/U rationale, and post-cutoff exclusion flag. The second is an \emph{institutional eligibility ledger}: every candidate capable of altering the oldest-four boundary, regulator identity, founding year, stratum and exclusion or tie reason. The third is the \emph{57-row outcome ledger}: final selected institution ID and classification, relevant local instrument, official direct policy locator when available, search surfaces, negative-search closure date, and rationale for $\stP$ or $\stU$. A script may then reconstruct the six $c_j$ vectors, $D$, the institutional margins, and each sensitivity without manual edits.

Publication in a repository must preserve the difference between a file that was actually retrieved and a link that is currently discoverable. The provenance manifest should separately identify the public URL, the archived local digest, and whether redistribution of the content is permitted. When only metadata/link sharing is authorized, the source can be cited and hashed without republishing the publisher's PDF. Hashes should bind the record and exact edition; a successful recomputation of the arithmetic is not independent verification of the underlying classification. The accompanying CSV links all 57 institutional labels to the locators \emph{already embedded in the submitted six-page paper}; the existence of those links alone does not close the 36-record source-extraction or second-coder requirement.

An independent auditor should freeze and time-stamp all source objects before seeing the primary table, then classify them in randomized order, blind to the target jurisdiction's headline result. For any conflict, the adjudication log should cite the controlling clause and explain whether the disagreement concerns authority, scope, document status, timing, or language. Raw pre-adjudication counts are indispensable: an agreed-upon final label may conceal poor initial reproducibility. Agreement should be reported on the actual number of independently coded objects, not on duplicated translations or repeated records, and separately for layer admission and institutional-panel coding when their units differ. A revised DOI-bearing public release would then make the protocol citable and independently inspectable, rather than using the word ``reproducible'' as a substitute for a released ledger.
